% Part: turing-machines 
% Chapter: undecidability 
% Section: verification

\documentclass[../../../include/open-logic-section]{subfiles}

\begin{document}

\olfileid{tur}{und}{ver} 
\olsection{نمائندگی دی تصدیق}

\begin{explain}
ساڈی نمائندگی دے درست کم کرن دی تصدیق لئی سانوں دو
گلاں ثابت کرنیاں پین گیاں۔ پہلے، ایہ دکھانا پئے گا
کہ جے~$M$ انپٹ~$w$ اُتے رُکدی اے، تاں
$!T(M, w) \lif !E(M, w)$ درست اے۔ فیر ایہدا اُلٹ
دکھانا پئے گا: جے~$!T(M, w) \lif !E(M, w)$ درست
اے، تاں~$M$ انپٹ~$w$ اُتے واقعی آخرکار
رُکدی اے۔

ایہ دو نتیجے ثابت کرن دیاں تدبیراں بہت وکھریاں نیں۔
پہلے نتیجے لئی سانوں دکھانا اے کہ پہلے درجے دی منطق
دا اک !!a{sentence}، یعنی~$!T(M, w) \lif !E(M, w)$،
درست اے۔ ایہ دکھان دا سب توں سوکھا طریقہ اک
!!a{derivation} دینا اے۔ پر ساڈا ثبوت ہر~$M$ تے~$w$
لئی ٹھیک ہووے؛ ایس لئی ایتھے لامتناہی جملے نیں۔ جیہڑا وی
!!{sentence} ہووے، اوہدا اک !!a{derivation} دینا پئے گا۔
ایس لئی اسی استقرا نال ایہ
ثابت کر سکدے آں کہ~$M$ تے~$w$ جیہو جیہے وی ہون،
تے~$M$ نوں انپٹ~$w$ اُتے رُکّن لئی جنے وی قدم
لگّن، $!T(M, w) \lif !E(M, w)$ دا اک
!!a{derivation} موجود ہووے گا۔

فطری طور اُتے ساڈا استقرا اُنہاں قدماں دی گنتی اُتے
ہووے گا جیہڑے~$M$ رُکّن والی مکمل حالت تک پہنچن
توں پہلاں چُکدی اے۔ استقرائی ثبوت وچ اسی قائم
کراں گے کہ ہر قدم~$n$ لئی، جیہڑا~$M$ انپٹ~$w$
اُتے چلن وچ چُکدی اے، $!T(M, w) \Entails !C(M, w, n)$،
جتھے~$!C(M, w, n)$، $M$ دے انپٹ~$w$ اُتے
چلدیاں~$n$ قدم چُکن توں بعد دی مکمل حالت ٹھیک
بیان کردا اے۔ ہن جے~$M$ انپٹ~$w$ اُتے، مثلاً،
$n$ قدم چُکن توں بعد رُکدی اے، تاں~$!C(M, w, n)$
رُکّن والی مکمل حالت بیان کرے گا۔ اسی ایہ وی
دکھاواں گے کہ~$!C(M, w, n)
\Entails !E(M, w)$ جدوں وی~$!C(M, w, n)$ ایسی رُکّن
والی مکمل حالت بیان کردا اے۔ ایس لئی جے~$M$ انپٹ~$w$
اُتے رُکدی اے، تاں کسے~$n$ لئی~$M$، $n$ قدم
چُکن توں بعد رُکّن والی مکمل حالت وچ ہووے گی۔
فیر~$!T(M, w) \Entails !C(M, w, n)$، جتھے
$!C(M, w, n)$ رُکّن والی مکمل حالت بیان کردا اے؛
تے ایس صورت وچ~$!C(M, w, n) \Entails !E(M, w)$
ہون کرکے سانوں~$T(M, w) \Entails !E(M, w)$ ملدا
اے، یعنی~$\Entails !T(M, w) \lif !E(M, w)$۔

اُلٹے رخ دی تدبیر بالکل وکھری اے۔ ایتھے اسی
$\Entails !T(M, w) \lif !E(M, w)$ فرض کرکے
ثابت کردے آں کہ~$M$ انپٹ~$w$ اُتے رُکدی اے۔
مفروضے توں~$!T(M, w) \Entails !E(M,
w)$ ملدا اے، یعنی ہر اُس !!{structure} وچ~$!E(M,
w)$ سچ اے جتھے~$!T(M, w)$ سچ اے۔ ایس
لئی اسی اک !!a{structure}~$\Struct{M}$ بیان کراں
گے جتھے~$!T(M, w)$ سچ اے: اوہدا دائرہ~$\Nat$
ہووے گا، تے سارے~$\Obj Q_q$ تے~$\Obj S_\sigma$
دیاں تعبیراں~$M$ دے انپٹ~$w$ اُتے چلن دیاں
مکمل حالتاں توں ملن گیاں۔ مثال لئی،
$\Sat{M}{\Obj Q_q(\num{m}, \num{n})}$ اُدوں تے صرف اُدوں
% PNB-CMP-031: source names T here; the running machine is M.
جدوں~$M$ انپٹ~$w$ اُتے~$n$ قدم چلن توں بعد
حالت~$q$ وچ اے تے خانہ~$m$ پڑھ رہی اے۔
ہن چونکہ مفروضے مطابق~$!T(M, w) \Entails !E(M, w)$،
تے تعمیر مطابق~$\Sat{M}{!T(M,
  w)}$، ایس لئی~$\Sat{M}{!E(M, w)}$۔ پر
$\Sat{M}{!E(M, w)}$ اُدوں تے صرف اُدوں جدوں کوئی
$n \in \Domain{M} = \Nat$ ایہو جیہا ہووے کہ~$M$
انپٹ~$w$ اُتے~$n$ قدم چُکن توں بعد رُکّن
والی مکمل حالت وچ ہووے۔
\end{explain}

\begin{defn} 
$!C(M, w, n)$ ایہ !!{sentence} ہووے:
\[ 
\Obj Q_q(\num{m}, \num{n}) \land \Obj S_{\sigma_0}(\num{0}, \num{n})
\land \dots \land \Obj S_{\sigma_k}(\num{k}, \num{n}) \land
\lforall[x][(\num{k} < x \lif \Obj S_\TMblank(x, \num{n}))]
\] 
جتھے~$q$، $M$ دی ویلے~$n$ والی حالت اے،
$M$ خانہ~$m$ ویلے~$n$ اُتے پڑھ رہی اے، خانے~$i$
وچ نشان~$\sigma_i$ ویلے~$n$ اُتے اے، $0
\le i \le k$ لئی؛ تے~$k$ یا تاں ویلے~$0$
اُتے پٹی دے سب توں سجّے غیر خالی خانے دا نمبر
اے یا~$n$ قدماں تک سر جیہڑے خانیاں تک پہنچ چکیا
اے، اوہناں وچ سب توں سجّے خانے دا نمبر، جیہڑا وی وڈا ہووے۔
\end{defn}

\begin{lem}\ollabel{lem:halt-config-implies-halt}
جے~$M$ انپٹ~$w$ اُتے~$n$ قدم چُکن
توں بعد رُکّن والی مکمل حالت وچ اے، تاں
$!C(M, w, n) \Entails !E(M, w)$۔
\end{lem}

\begin{proof}
فرض کرو~$M$ انپٹ~$w$ اُتے~$n$ قدم چُکن توں
بعد رُکدی اے۔ کوئی حالت~$q$، خانہ~$m$ تے نشان
~$\sigma$ ایداں دے نیں کہ:
\begin{enumerate} 
\item $n$ قدم چُکن توں بعد~$M$ حالت~$q$ وچ
  خانہ~$m$ پڑھ رہی اے، جتھے~$\sigma$ لکھیا اے۔
\item حالتی تبدیلی دا فنکشن~$\delta(q, \sigma)$
  غیر تعریف شدہ اے۔
\end{enumerate}
$!C(M, w, n)$ ایس مکمل حالت دا بیان اے تے ایس
وچ اجزا~$\Obj Q_{q}(\num{m}, \num{n})$ تے~$\Obj
S_{\sigma}(\num{m}, \num{n})$ شامل ہون گے۔ ایہ
اجزا مل کے~$!E(M,
w)$ لازم کردے نیں:
\[
\lexists[x][\lexists[y][(\bigvee_{\tuple{q, \sigma} \in
      X}(\Obj Q_q(x, y) \land \Obj S_\sigma(x, y)))]]
\]
% PNB-CMP-032: use the stated halting pair and the S predicate symbol.
کیوں جے~$\Obj Q_{q}(\num{m}, \num{n}) \land \Obj S_{\sigma}(\num{m},
\num{n}) \Entails \bigvee_{\tuple{q, \sigma} \in X} (\Obj Q_q(\num{m},
\num{n}) \land \Obj S_{\sigma}(\num{m}, \num{n}))$،
ایس لئی کہ~$\tuple{q,\sigma} \in X$ اے۔
\end{proof}

\begin{explain}
ایس لئی جے~$M$ انپٹ $w$ اُتے رُکدی اے، تاں
کوئی~$n$ ایہو جیہا اے کہ~$!C(M,
w, n) \Entails !E(M,w)$۔ ہن اسی دکھاواں گے
کہ ہر اُس ویلے~$n$ لئی، جتھے مشین پہلاں نہ رُکی
ہووے، $!T(M, w) \Entails !C(M, w, n)$۔
\end{explain}

\begin{lem}
\ollabel{lem:config}
% PNB-CMP-033: include the configuration at the first halting step.
ہر~$n$ لئی، جے~$M$، $n$ قدم چُکن توں پہلاں
نہ رُکی ہووے، تاں~$!T(M, w)
\Entails !C(M, w, n)$۔
\end{lem}

\begin{proof}
استقرا دی بنیاد: جے~$n = 0$، تاں~$!C(M, w, 0)$
دے عطف دے سارے جزو~$!T(M, w)$ دے عطف وچ وی نیں۔ ایس
لئی اوہ ایس توں منطقی طور اُتے نکلدے نیں۔

استقرائی مفروضہ: جے~$M$، $n$ویں قدم توں پہلاں
نہ رُکی ہووے، تاں~$!T(M,w) \Entails !C(M, w, n)$۔
سانوں دکھانا اے کہ (جے~$!C(M, w, n)$ رُکّن والی
مکمل حالت بیان نہ کردا ہووے) تاں~$!T(M, w) \Entails
!C(M, w, n+1)$۔

% PNB-CMP-034: induction step also needs n=0.
فرض کرو~$n \ge 0$، تے~$n$ قدم چُکن توں بعد
$M$، انپٹ~$w$ اُتے شروع ہو کے، حالت~$q$ وچ
خانہ~$m$ پڑھ رہی اے۔ چونکہ~$M$، $n$ قدم چُکن
توں بعد نہیں رُکدی، ہِٹھلیاں تِن صورتاں وچوں
کسے اک صورت دا حکم~$M$ دے
پروگرام وچ موجود ہووے گا:

\begin{enumerate} 
\item \ollabel{right} $\delta(q, \sigma) = \tuple{q', \sigma', \TMright}$

\item \ollabel{left} $\delta(q, \sigma) = \tuple{q', \sigma', \TMleft}$

\item \ollabel{stay} $\delta(q, \sigma) = \tuple{q', \sigma', \TMstay}$
\end{enumerate}

اسی ایہناں تِن صورتاں نوں واری واری ویکھاں گے۔

\begin{enumerate} 
\item فرض کرو~\olref{right} والے روپ دا حکم موجود اے۔
  \olref[rep]{defn:tm-descr}\olref[rep]{rep-right} مطابق ایہدے
  معنی ایہ نیں کہ
% PNB-CMP-035: omit the old scanned square m from retained-symbol conjuncts.
\begin{align*} 
& \lforall[x][\lforall[y][((\Obj Q_{q}(x,
  y) \land \Obj S_\sigma(x, y)) \lif {}]]\\
& \qquad (\Obj Q_{q'}(x',
  y') \land \Obj S_{\sigma'}(x, y') \land !A(x, y)))  
\intertext{، $!T(M,w)$ دا اک جزو اے۔ ایس توں ہِٹھلا
  !!{sentence} نکلدا اے (کُلی تخصیص نال~$\num{m}$ نوں
  $x$ دی تھاں تے~$\num{n}$ نوں~$y$ دی تھاں رکھ کے):}
& (\Obj Q_{q}(\num{m}, \num{n}) \land \Obj S_{\sigma}(\num{m},
\num{n})) \lif {}\\
& \qquad (\Obj Q_{q'}(\num{m}', \num{n}') \land
\Obj S_{\sigma'}(\num{m}, \num{n}') \land !A(\num{m}, \num{n})).
\intertext{استقرائی مفروضے مطابق~$!T(M, w) \Entails !C(M, w, n)$،
  یعنی:}
& \Obj Q_q(\num{m}, \num{n}) \land \Obj S_{\sigma_0}(\num{0}, \num{n})
\land \dots \land \Obj S_{\sigma_k}(\num{k}, \num{n}) \land \\
& \qquad\lforall[x][(\num{k} < x \lif \Obj S_\TMblank(x, \num{n}))]\\
\intertext{چونکہ~$n$ قدم چُکن توں بعد خانہ~$m$ وچ~$\sigma$
  اے، ایس لئی متعلقہ جزو~$\Obj S_\sigma(\num{m}, \num{n})$
  اے؛ لہٰذا ایہ لازم آندا اے:}
& \Obj Q_{q}(\num{m}, \num{n}) \land \Obj S_{\sigma}(\num{m},
   \num{n})
\intertext{ہن سانوں ایہ ملدا اے:}
& \Obj Q_{q'}(\num{m}', \num{n}') \land \Obj S_{\sigma'}(\num{m},
  \num{n}') \land {}\\
& \qquad\bigwedge_{\substack{0\le i\le k\\ i\ne m}}
  \Obj S_{\sigma_i}(\num{i}, \num{n}') \land {}\\
& \qquad \lforall[x][(\num{k} < x \lif \Obj S_\TMblank(x, \num{n}'))]
\end{align*}
ایہ نتیجہ ایویں ملدا اے: پہلی لکیر پچھلے شرطیہ دے
تالی توں اثباتِ مقدم نال سدھی نکلدی اے۔ وچلی
لکیر وچ~$S_{\sigma_m}(\num{m},\num{n}')$ شامل نہیں۔
ایس لکیر دا ہر جزو $!C(M, w, n)$ دے متعلقہ جزو تے
$!A(\num{m}, \num{n})$ توں نکلدا اے۔

جے~$m < k$، تاں~$!T(M,w) \Proves \num{m} < \num {k}$
(\olref[rep]{prop:mlessk})، تے~$<$ دی تعدیت توں
سانوں~$\lforall[x][(\num{k} < x \lif \num{m} < x)]$
ملدا اے۔ جے~$m = k$، تاں ایہو
$\lforall[x][(\num{k} < x \lif \num{m} < x)]$
صرف منطق نال نکلدا اے۔ آخری لکیر فیر
$!C(M, w,
n)$ دے متعلقہ جزو، $\lforall[x][(\num{k} < x \lif \num{m} < x)]$،
تے~$!A(\num{m},
\num{n})$ توں نکلدی اے۔ جے~$m<k$، تاں ایہ
پہلاں ہی~$!C(M, w, n+1)$ اے۔

ہن فرض کرو~$m=k$۔ ایس صورت وچ~$n+1$ قدم
چُکن توں بعد پٹی دا سر خانہ~$k+1$ تک وی پہنچ چکیا
اے۔ ہن سر جیہڑے خانیاں تک پہنچیا اے، اوہناں وچوں ایہ سب توں سجّا اے۔
ایس لئی~$!C(M, w, n+1)$ وچ اک نواں جزو
$\Obj S_\TMblank(\num{k}',\num{n}')$ اے، تے آخری جزو
$\lforall[x][(\num{k}' < x \lif \Obj S_\TMblank(x, \num{n}'))]$
اے۔ سانوں جانچنا اے کہ ایہ دونیں !!{sentence}s
وی منطقی طور اُتے نکلدے نیں۔

ساڈے کول پہلے ہی
$\lforall[x][(\num{k} < x \lif \Obj S_\TMblank(x,
  \num{n}'))]$ اے۔ خاص طور اُتے ایہدے توں
$\num{k} < \num{k}' \lif
\Obj S_\TMblank(\num{k}', \num{n}')$ ملدا اے۔ مسلّمے
$\lforall[x][x <
  x']$ توں~$\num{k} < \num{k}'$ ملدا اے۔ اثباتِ
مقدم نال~$\Obj
S_\TMblank(\num{k}',\num{n}')$ نکلدا اے۔

نالے، چونکہ~$!T(M,w) \Proves \num{k} < \num{k}'$،
تاں~$<$ دی تعدیت دا مسلّمہ سانوں
$\lforall[x][(\num{k}' < x \lif \Obj
  S_\TMblank(x, \num{n}'))]$ دیندا اے۔
(ایس دی جانچ اسی مشق لئی چھڈدے آں۔)

\item فرض کرو~\olref{left} والے روپ دا حکم موجود اے۔
  فیر~\olref[rep]{defn:tm-descr}\olref[rep]{rep-left} مطابق:
% PNB-CMP-035: omit the written square from the left-case retained-symbol conjunction.
% PNB-CMP-036: left-frame exemption is the written x-prime; edge states match q and q-prime.
\begin{align*} 
& \lforall[x][\lforall[y][((\Obj Q_{q}(x', y) \land \Obj
    S_{\sigma}(x', y)) \lif {}]]\\
& \qquad (\Obj Q_{q'}(x, y') \land \Obj
  S_{\sigma'}(x', y') \land !A(x', y))) \land {}\\
& \lforall[y][((\Obj Q_{q}(\num{0}, y) \land \Obj S_{\sigma}(\num{0},
    y)) \lif {}]\\
& \qquad (\Obj Q_{q'}(\num{0}, y') \land \Obj S_{\sigma'}(\num{0}, y')
  \land !A(\num{0}, y)))
\intertext{، $!T(M,w)$ دا اک جزو اے۔ جے~$m>0$،
  تاں~$l = m - 1$ رکھو (یعنی~$m = l+1$)۔ اُتلے
  !!{sentence} دے پہلے جزو توں ایہ نتیجہ نکلدا اے:}
& (\Obj Q_{q}(\num{l}', \num{n}) \land \Obj S_{\sigma}(\num{l}', \num{n}))
\lif {} \\
& \qquad (\Obj Q_{q'}(\num{l}, \num{n}') \land \Obj S_{\sigma'}(\num{l}',
\num{n}') \land !A(\num{l}', \num{n}))
\intertext{دوجی صورت وچ~$l = m = 0$ رکھو تے دوجے
  جزو توں نکلن والا ہِٹھلا !!{sentence} ویکھو:}
& ((\Obj Q_{q}(\num{0}, \num{n}) \land \Obj S_{\sigma}(\num{0}, \num{n})) \lif {}\\
& \qquad   (\Obj Q_{q'}(\num{0}, \num{n}') \land \Obj S_{\sigma'}(\num{0}, \num{n}') \land
!A(\num{0}, \num{n}))) 
\intertext{دوہاں وچوں ہر اک جملے نال ایہ نتیجہ نکلدا اے:}
&  \Obj Q_{q'}(\num{l}, \num{n}') \land \Obj S_{\sigma'}(\num{m},
  \num{n}') \land {}\\
&\qquad  \bigwedge_{\substack{0\le i\le k\\ i\ne m}}
  \Obj S_{\sigma_i}(\num{i}, \num{n}') \land {} \\
& \qquad  \lforall[x][(\num{k} < x
  \lif \Obj S_\TMblank(x, \num{n}'))]
\end{align*}
پہلاں وانگ۔ (دھیان دیو کہ پہلی صورت وچ
$\num{l}' \ident \num{l+1}
\ident \num{m}$ تے دوجی صورت وچ
$\num{l} \ident \num{0}$۔) پر ایہ عین
$!C(M, w, n+1)$ اے۔

\item صورت~\olref{stay} مشق لئی چھڈدے آں۔
\end{enumerate}
اسی دکھا دتا کہ ہر اُس~$n$ لئی، جیہدے توں پہلاں
مشین نہ رُکی ہووے، $!T(M, w) \Entails !C(M, w, n)$۔
\end{proof}

\begin{prob}
\olref[tur][und][ver]{lem:config} دے ثبوت وچ
صورت~\olref[tur][und][ver]{stay} مکمل کرو۔
\end{prob}

\begin{prob}
$\Obj S_{\sigma_i}(\num{i}, \num{n}')$ دا اک
!!a{derivation}، $\Obj S_{\sigma_i}(\num{i}, \num{n})$
تے~$!A(m, n)$ توں دیو (ایس مفروضے ہِٹھ کہ~$i \neq
m$، یعنی یا~$i < m$ یا~$m < i$)۔
\end{prob}

\begin{prob}
$\lforall[x][(\num{k}' < x \lif \Obj
  S_\TMblank(x, \num{n}'))]$ دا اک !!a{derivation}
$\lforall[x][(\num{k} < x \lif \Obj
  S_\TMblank(x, \num{n}'))]$، $\lforall[x][x < x']$ تے
$\lforall[x][\lforall[y][\lforall[z][ ((x < y \land y < z) \lif x <
      z)]]]$ توں دیو۔)
\end{prob}


\begin{lem}
\ollabel{lem:valid-if-halt}
جے~$M$ انپٹ~$w$ اُتے رُکدی اے، تاں
$!T(M, w) \lif
!E(M, w)$ درست اے۔
\end{lem}

\begin{proof}
\olref{lem:config} توں سانوں پتا اے کہ رُکّن توں
پہلاں یا رُکّن دے ویلے~$n$
اُتے مکمل حالت دا بیان~$!C(M, w, n)$، جیہڑا~$M$
دی ویلے~$n$ اُتے مکمل حالت دسدے، $!T(M, w)$
توں نکلدا اے۔ فرض کرو~$M$، $k$ قدم چُکن توں
بعد رُکدی اے۔ اوس ویلے اوہ خانہ~$m$ پڑھ رہی
ہووے گی، کسے~$m \in \Nat$ لئی۔ فیر
$!C(M, w, k)$، $M$ دی رُکّن والی مکمل حالت
بیان کردا اے؛ یعنی ایس وچ~$\Obj Q_q(\num{m}, \num{k})$
تے~$\Obj
S_\sigma(\num{m}, \num{k})$ دونیں جزو شامل نیں،
جدوکہ~$\delta(q,\sigma)$ غیر تعریف شدہ اے۔
ایس لئی~\olref{lem:halt-config-implies-halt} مطابق
$!C(M, w, k) \Entails !E(M,
w)$۔ پر چونکہ~$!T(M, w) \Entails !C(M, w, k)$،
ایس لئی~$!T(M, w)
\Entails !E(M, w)$، تے نتیجتاً
$!T(M, w) \lif !E(M, w)$ درست اے۔
\end{proof}


\begin{explain} 
اپنے دعوے دی تصدیق مکمل کرن لئی سانوں اُلٹا
رخ وی قائم کرنا پئے گا: جے~$!T(M, w) \lif !E(M, w)$
درست اے، تاں~$M$ انپٹ~$w$ اُتے واقعی
رُکدی اے۔
\end{explain}

\begin{lem}
\ollabel{lem:halt-if-valid}
جے~$\Entails !T(M, w) \lif !E(M, w)$، تاں~$M$
انپٹ~$w$ اُتے رُکدی اے۔
\end{lem}

\begin{proof}
$\Lang L_M$ دی !!{structure}~$\Struct M$ ویکھو
جیہدا دائرہ~$\Nat$ اے۔ ایہ~$\Obj 0$ دی تعبیر~$0$،
$\prime$ دی تعبیر جانشین فنکشن تے~$<$ دی تعبیر “توں گھٹ”
والے تعلق دے طور اُتے کردی اے؛ تے محمولاں~$\Obj Q_q$
تے~$\Obj S_\sigma$ دی ایویں:
\begin{align*}
  \Assign{\Obj Q_q}{M} & =
\Setabs{\tuple{m, n}}{\begin{array}{ll}\text{انپٹ~$w$ اُتے شروع ہو کے، $n$ قدم بعد،}\\ \text{$M$ حالت $q$
  وچ خانہ~$m$ پڑھ رہی اے}\end{array}} \\
\Assign{\Obj S_\sigma}{M} & = \Setabs{\tuple{m, n}}{\begin{array}{ll}
\text{انپٹ~$w$ اُتے شروع ہو کے، $n$ قدم بعد،}\\ \text{خانہ~$m$، $M$ دا، وچ
  نشان~$\sigma$ اے}\end{array}}
\end{align*}
ہور لفظاں وچ، اسی !!{structure}~$\Struct{M}$
ایویں بناندے آں کہ اوہ~$M$ دا انپٹ~$w$
اُتے حقیقی چلن قدم بہ قدم بیان کرے۔
% PNB-CMP-037: extend a finite run by retaining its final halting configuration.
جے مشین رُک جائے تاں بعدلے ویلّیاں لئی آخری
رُکّن والی مکمل حالت اوہو جیہی رکھو؛ ایویں
تعبیراں سارے قدرتی عددی ویلّیاں اُتے مقرر ہو جاندیاں نیں۔
صاف اے کہ~$\Sat{M}{!T(M, w)}$۔ جے
$\Entails !T(M, w) \lif !E(M, w)$،
تاں~$\Sat{M}{!E(M, w)}$ وی، یعنی:
\[
\Sat{M}{\lexists[x][\lexists[y][(\bigvee_{\tuple{q, \sigma} \in
      X}(\Obj Q_q(x, y) \land \Obj S_\sigma(x, y)))]]}.
\]
چونکہ~$\Domain{M} = \Nat$، ایس لئی کوئی
$m$، $n \in \Nat$ ضرور ہون گے جنہاں لئی
$\Sat{M}{\Obj Q_q(\num{m}, \num{n}) \land \Obj S_\sigma(\num{m},
\num{n})}$، کسے~$q$ تے~$\sigma$ لئی، جتھے~$\delta(q, \sigma)$ غیر تعریف
شدہ اے۔ $\Struct M$ دی تعریف مطابق ایہدے معنی
ایہ نیں کہ~$M$، انپٹ~$w$ اُتے شروع ہو کے، $n$
قدم چُکن توں بعد حالت~$q$ وچ اے تے نشان~$\sigma$
پڑھ رہی اے، تے حالتی تبدیلی دا فنکشن غیر تعریف
شدہ اے؛ یعنی~$M$ رُک چکی اے۔
\end{proof}

\end{document}
